% ============================================================
% FEE3411 Control Systems A -- Tutorial sheet, Week 02
% Topic: Electromechanical modelling and linear approximation
%
% ONE source, TWO outputs. The solutions toggle is driven by \jobname,
% so the two PDFs cannot drift apart:
%
%   latexmk -pdf -jobname=week-02-tutorial-questions week-02-tutorial.tex
%   latexmk -pdf -jobname=week-02-tutorial-solutions week-02-tutorial.tex
%
% Needs: ../tex/preamble.tex (this file lives in Tutorials/).
% ============================================================
\documentclass[11pt,a4paper]{article}
\usepackage[margin=25mm]{geometry}
% ============================================================
% FEE3411 Control Systems A -- shared preamble
% Used by both the lecture notes (article) and the slides (beamer).
% ============================================================
\usepackage[T1]{fontenc}
% lmodern gives cleaner PDF fonts but is absent from minimal TeX installs.
% Load it only if present, so this preamble builds anywhere.
\IfFileExists{lmodern.sty}{\usepackage{lmodern}}{}
\usepackage{amsmath,amssymb,amsthm}
\usepackage{graphicx}
\usepackage{booktabs}
\usepackage{siunitx}
\usepackage{tikz}
\usepackage{circuitikz}
\usepackage{pgfplots}
\pgfplotsset{compat=1.18}
\usetikzlibrary{arrows.meta,positioning,calc,decorations.markings,shapes.geometric}
\usepackage[most]{tcolorbox}
\usepackage{xcolor}

% ---- Course colours -------------------------------------------------
\definecolor{fnavy}{HTML}{1F3864}
\definecolor{faccent}{HTML}{2E5E4E}
\definecolor{fflag}{HTML}{9C4221}

% ---- Block-diagram style --------------------------------------------
\tikzset{
  block/.style   = {draw, thick, rectangle, minimum height=2.6em,
                    minimum width=3.2em, align=center},
  sumnode/.style = {draw, thick, circle, inner sep=0pt, minimum size=1.6em},
  pickoff/.style = {fill, circle, inner sep=0pt, minimum size=3pt},
  sig/.style     = {-{Latex[length=2mm]}, thick},
}

% ---- Summing-junction signs -----------------------------------------
% \sumsign{<node>}{<angle>}{<+ or ->} places a sign just outside the circle.
\newcommand{\sumsign}[3]{\node[font=\small] at ($(#1)+(#2:0.5)$) {$#3$};}

% ---- Boxed environments ---------------------------------------------
\newtcolorbox{keyidea}[1][]{colback=fnavy!5,colframe=fnavy,
  fonttitle=\bfseries,title=Key idea,#1}
\newtcolorbox{workedex}[1][]{colback=faccent!5,colframe=faccent,
  fonttitle=\bfseries,title=Worked example,#1}
\newtcolorbox{pitfall}[1][]{colback=fflag!5,colframe=fflag,
  fonttitle=\bfseries,title=Common pitfall,#1}

% ---- Shorthands ------------------------------------------------------
\newcommand{\Lap}{\mathcal{L}}
\newcommand{\jw}{\mathrm{j}\omega}
\newcommand{\wn}{\omega_{n}}
\newcommand{\zt}{\zeta}
\newcommand{\TF}[1]{#1(s)}

% ---- Week title machinery -------------------------------------------
\usepackage{fancyhdr}
\newcommand{\thecoursecode}{FEE3411}
\newcommand{\thecoursename}{Control Systems A}
\newcommand{\wknum}{}
\newcommand{\wktopic}{}
\newcommand{\weeknumber}[1]{\renewcommand{\wknum}{#1}}
\newcommand{\weektopic}[1]{\renewcommand{\wktopic}{#1}}
\newcommand{\weektitle}{%
  \thispagestyle{plain}%
  \begin{center}
    {\large\sffamily\bfseries\color{fnavy}\thecoursecode: \thecoursename}\\[2pt]
    {\Huge\sffamily\bfseries Week \wknum}\\[4pt]
    {\LARGE\sffamily\color{fnavy}\wktopic}\\[6pt]
    \rule{\textwidth}{1.2pt}
  \end{center}\vspace{4pt}%
  \pagestyle{fancy}\fancyhf{}%
  \fancyhead[L]{\footnotesize\color{gray}\thecoursecode\ Week \wknum\ --- \wktopic}%
  \fancyhead[R]{\footnotesize\color{gray}\thepage}%
  \renewcommand{\headrulewidth}{0.4pt}%
}

% ---- Outcomes and reading boxes -------------------------------------
\newenvironment{outcomes}
  {\begin{tcolorbox}[colback=white,colframe=fnavy!60,
      fonttitle=\bfseries,title=By the end of this week you should be able to]
   \begin{itemize}[leftmargin=*]}
  {\end{itemize}\end{tcolorbox}}

\newtcolorbox{readingbox}{colback=gray!5,colframe=gray!50,
  fonttitle=\bfseries,title=Reading}

% enumitem must NOT be loaded under beamer: it overwrites beamer's itemize
% templates and the bullet markers silently disappear. The notes (article) get
% it; the slides (beamer) use beamer's own list machinery instead.
\makeatletter
\@ifclassloaded{beamer}{}{\usepackage{enumitem}}
\makeatother

\usepackage{environ}

% ---- solutions toggle, decided by the job name ----------------------
% NB: xstring's \IfSubStr does NOT work here -- \jobname yields catcode-12
% characters and the comparison silently fails, giving the questions version
% both times. The expl3 string test below is catcode-blind and does work.
\newif\ifsolutions
\ExplSyntaxOn
\tl_set:Nx \l_tmpa_tl { \jobname }
\str_if_in:NnTF \l_tmpa_tl { solutions } { \solutionstrue } { \solutionsfalse }
\ExplSyntaxOff

\newtcolorbox{solbox}[1][]{colback=faccent!4,colframe=faccent!70,
  fonttitle=\bfseries,title=Solution,breakable,#1}
\NewEnviron{solution}{\ifsolutions\begin{solbox}\BODY\end{solbox}\fi}

% ---- question tagging ------------------------------------------------
\newcommand{\inhour}{\hfill{\small\color{faccent}\itshape[in the hour]}}
\newcommand{\athome}{\hfill{\small\color{fflag}\itshape[homework]}}

\weeknumber{2}
\weektopic{Tutorial 2 --- Electromechanical Modelling and Linearization}

\begin{document}
\weektitle

\begin{center}
\ifsolutions
  {\large\sffamily\bfseries\color{fflag}QUESTIONS AND SOLUTIONS}\\[2pt]
  {\small\itshape Instructor copy --- do not circulate before the tutorial.}
\else
  {\large\sffamily\bfseries\color{fnavy}QUESTION SHEET}
\fi
\end{center}

\bigskip
\noindent
This hour has two purposes: turning the servomotor derivation from the lecture
into working fluency with numbers, and extending gearing and linearization
slightly beyond what the notes worked through, so you meet the ideas in a
second setting before Assignment~1 asks you to use them in a third.

\smallskip
\noindent
Questions marked \textbf{\color{faccent}[in the hour]} are the ones we work
through together. Those marked \textbf{\color{fflag}[homework]} you do in your
own time --- they are not collected, but the material is assumed from here on,
including in Assignment~1 (Questions~3--4, due Week~3).

\smallskip
\noindent
\emph{Preparation:} Week~2 lecture notes, \S6--\S7.

% --------------------------------------------------------------------
\begin{readingbox}
\textbf{Quick reference --- carry these into every question below.}

\smallskip
\begin{minipage}[t]{0.48\textwidth}
\textbf{Armature-controlled d.c.\ servomotor}\\[2pt]
\footnotesize
\[
  \frac{\theta_m(s)}{V_a(s)}=\frac{K_t}{s\big[(R_a+L_as)(J_ms+b)+K_tK_b\big]}
\]
\[
  L_a\to0:\quad \frac{\theta_m(s)}{V_a(s)}=\frac{K_m}{s(\tau_ms+1)},
\]
\[
  K_m=\frac{K_t}{R_ab+K_tK_b},\qquad \tau_m=\frac{J_mR_a}{R_ab+K_tK_b}
\]
\end{minipage}
\hfill
\begin{minipage}[t]{0.48\textwidth}
\textbf{Gear reflection to the motor shaft}\\[2pt]
\footnotesize
\[
  n=\frac{\theta_L}{\theta_m},\qquad
  J_{\text{eq}}=J_m+n^2J_L,\qquad
  b_{\text{eq}}=b+n^2b_L
\]
\medskip
\textbf{Linearization about $(x_0,u_0)$}
\[
  \Delta\dot x \approx
  \left.\frac{\partial f}{\partial x}\right|_{0}\!\Delta x
  +\left.\frac{\partial f}{\partial u}\right|_{0}\!\Delta u,
  \qquad f(x_0,u_0)=0
\]
\end{minipage}
\end{readingbox}

% ====================================================================
\section*{Question 1 --- The servomotor: derive, reduce, and check the
reduction \inhour}

An armature-controlled d.c.\ servomotor has $R_a=\SI{4}{\ohm}$,
$L_a=\SI{0.05}{\henry}$, $K_t=\SI{0.8}{\newton\metre\per\ampere}$,
$K_b=\SI{0.8}{\volt\second\per\radian}$,
$J_m=\SI{0.05}{\kilogram\metre\squared}$ and
$b=\SI{0.06}{\newton\metre\second\per\radian}$. The field is separately
excited and held constant.

\begin{enumerate}[label=(\alph*)]
\item Write down the armature circuit equation and the torque-balance
      equation, and state the transfer function $\theta_m(s)/V_a(s)$ that
      follows from eliminating the armature current. (This is bookwork from
      the lecture --- quote the result, but be ready to reproduce the
      elimination step if asked.)

\item Setting $L_a\to0$, find the reduced transfer function
      $K_m/[s(\tau_ms+1)]$ numerically.

\item Without discarding $L_a$, find the two roots of the quadratic factor
      $(R_a+L_as)(J_ms+b)+K_tK_b=0$. Identify the \emph{dominant} pole (the
      one closest to the origin) and its time constant.

\item Compare the time constant from (c) with $\tau_m$ from (b). By what
      percentage do they differ? Was neglecting $L_a$ reasonable here?
\end{enumerate}

\begin{solution}
\textbf{(a)} Exactly as derived in the lecture (notes \S6.1--6.2):
\[
  v_a=R_ai_a+L_a\frac{di_a}{dt}+K_b\dot\theta_m,
  \qquad
  K_ti_a=J_m\ddot\theta_m+b\dot\theta_m,
\]
\[
  \frac{\theta_m(s)}{V_a(s)}=\frac{K_t}{s\big[(R_a+L_as)(J_ms+b)+K_tK_b\big]}.
\]

\medskip
\textbf{(b)} $R_ab+K_tK_b=(4)(0.06)+(0.8)(0.8)=0.24+0.64=0.88$.
\[
  K_m=\frac{0.8}{0.88}=\SI{0.909}{\radian\per\volt\per\second},
  \qquad
  \tau_m=\frac{(0.05)(4)}{0.88}=\SI{0.227}{\second}.
\]

\medskip
\textbf{(c)} Expanding $(R_a+L_as)(J_ms+b)+K_tK_b$ with the numbers above,
\[
  0.0025\,s^2+0.203\,s+0.88=0
  \qquad\Longleftrightarrow\qquad
  s^2+81.2s+352=0,
\]
with roots
\[
  s=-4.595 \quad\text{and}\quad s=-76.60 .
\]
The dominant pole is $s=-4.595$ (closest to the origin, so it decays slowest
and dominates the response), giving time constant
$\tau_{\text{dominant}}=1/4.595=\SI{0.218}{\second}$.

\medskip
\textbf{(d)} $\tau_m=\SI{0.227}{\second}$ against
$\tau_{\text{dominant}}=\SI{0.218}{\second}$ is a difference of about
\[
  \frac{0.227-0.218}{0.218}\times100\%\approx4.4\%.
\]
That is a small error for an introductory model, so neglecting $L_a$ was
reasonable here --- but notice it is not \emph{zero}. The fast pole at
$s=-76.6$ (time constant \SI{0.013}{\second}) is the one thrown away entirely;
it is legitimate to discard only because it is roughly seventeen times faster
than the dominant pole it leaves behind.
\end{solution}

\begin{pitfall}
A $4$--$5\%$ shift in the dominant time constant is the kind of number you
should expect to see and accept, not a sign of an arithmetic mistake. If the
same calculation ever gives you a discrepancy of, say, $50\%$, that is telling
you $L_a$ is \emph{not} negligible for that motor, and the reduced formula
should not be used --- check by comparing the electrical time constant
$L_a/R_a$ against the mechanical one, $\tau_m$; the approximation is only good
when the first is much smaller than the second.
\end{pitfall}

% ====================================================================
\section*{Question 2 --- A two-stage gear train \inhour}

The servomotor of Question~1 drives a load of inertia
$J_L=\SI{3}{\kilogram\metre\squared}$ through \emph{two} meshed gear pairs in
series: the first stage has ratio $n_1=\theta_1/\theta_m=1/2$ (an intermediate
shaft at angle $\theta_1$), and the second has ratio
$n_2=\theta_L/\theta_1=1/4$. Neglect the inertia of the intermediate shaft and
all load-side friction.

\begin{enumerate}[label=(\alph*)]
\item Find the overall ratio $n=\theta_L/\theta_m$.
\item Find the equivalent inertia $J_{\text{eq}}$ at the motor shaft.
\item Recompute $\tau_m$ for the loaded, geared motor, and compare it with the
      \emph{unloaded} motor's $\tau_m$ from Question~1(b). By what factor has
      the motor slowed down?
\end{enumerate}

\begin{solution}
\textbf{(a)} Ratios of gear pairs in series multiply:
\[
  n=n_1n_2=\left(\frac{1}{2}\right)\left(\frac{1}{4}\right)=\frac{1}{8}.
\]
This is the same idea as combining several transfer functions in a series
connection --- each stage's ratio is applied in turn.

\medskip
\textbf{(b)} Using $J_{\text{eq}}=J_m+n^2J_L$ with the \emph{overall} ratio
(the intermediate shaft's own inertia is neglected, so it never enters the
formula):
\[
  J_{\text{eq}}=0.05+\left(\frac{1}{8}\right)^2(3)=0.05+\frac{3}{64}
   =0.05+0.046875=\SI{0.096875}{\kilogram\metre\squared}.
\]

\medskip
\textbf{(c)}
\[
  \tau_m'=\frac{J_{\text{eq}}R_a}{R_ab+K_tK_b}=\frac{(0.096875)(4)}{0.88}
   =\SI{0.440}{\second}.
\]
Compared with the unloaded $\tau_m=\SI{0.227}{\second}$, the ratio is
\[
  \frac{0.440}{0.227}\approx1.94,
\]
so the geared, loaded motor is just under twice as sluggish as the bare motor
--- a modest penalty for driving a load $3/0.05=60$ times the motor's own
inertia, precisely because the $n^2=1/64$ scaling in the reflected inertia
does almost all of the work.
\end{solution}

% ====================================================================
\section*{Question 3 --- Linearizing a liquid-level tank \inhour}

A tank of constant cross-sectional area $A=\SI{2}{\metre\squared}$ has inflow
rate $q_{in}(t)$ (the control input) and an outflow through a valve at the
base obeying Torricelli's law, so the level $h(t)$ obeys
\[
  A\,\dot h(t) = q_{in}(t) - k\sqrt{h(t)}, \qquad k=\SI{0.5}{\metre^{1.5}\per\second}.
\]
This is nonlinear because of the $\sqrt{h}$ term.

\begin{enumerate}[label=(\alph*)]
\item Find the equilibrium inflow $q_{in,0}$ that holds the level at
      $h_0=\SI{4}{\metre}$.
\item Writing $\dot h=f(h,q_{in})$, find $\partial f/\partial h$ and
      $\partial f/\partial q_{in}$ at the operating point, and hence the
      linearized equation in $\Delta h$ and $\Delta q_{in}$.
\item Identify the time constant of the linearized tank, and state in words
      what it means physically for the tank to respond \emph{faster} or
      \emph{slower} than this.
\item The valve is replaced by a narrower one with the same $\sqrt{h}$ law but
      a smaller $k$. Without recomputing anything, say whether the tank's
      linearized time constant gets longer or shorter, and why.
\end{enumerate}

\begin{solution}
\textbf{(a)} At equilibrium $\dot h=0$, so $q_{in,0}=k\sqrt{h_0}=0.5\sqrt4=1$,
i.e.\ $q_{in,0}=\SI{1}{\metre\cubed\per\second}$.

\medskip
\textbf{(b)} With $f(h,q_{in})=(q_{in}-k\sqrt h)/A$,
\[
  \frac{\partial f}{\partial h}=-\frac{k}{2A\sqrt h},
  \qquad
  \frac{\partial f}{\partial q_{in}}=\frac{1}{A}.
\]
At $h_0=4$: $\partial f/\partial h=-0.5/(2\cdot2\cdot2)=-1/16=-0.0625$, and
$\partial f/\partial q_{in}=1/2=0.5$. So
\[
  \boxed{\;\Delta\dot h = 0.5\,\Delta q_{in} - 0.0625\,\Delta h\;} .
\]

\medskip
\textbf{(c)} Writing the linearized equation as
$\Delta\dot h+0.0625\,\Delta h=0.5\,\Delta q_{in}$, the time constant is
$\tau=1/0.0625=\SI{16}{\second}$. A tank with a \emph{shorter} time constant
would settle to a new level faster after a change in inflow; one with a
\emph{longer} time constant responds more sluggishly --- exactly the same
sense in which $\tau_m$ described the servomotor's speed of response.

\medskip
\textbf{(d)} A smaller $k$ makes $|\partial f/\partial h|=k/(2A\sqrt{h_0})$
\emph{smaller} in magnitude, so the time constant $\tau=2A\sqrt{h_0}/k$
\emph{increases} --- the tank becomes slower. Physically, a narrower valve
drains more slowly for a given head, so a disturbed level takes longer to
recover; this needs no recalculation because the sign and shape of the
dependence on $k$ are already visible in the linearized coefficient.
\end{solution}

\begin{keyidea}
The tank is the same kind of exercise as the pendulum in the lecture notes,
but with a control input as well as a state: linearizing
$\dot x=f(x,u)$ needs \emph{two} partial derivatives, one for each variable,
evaluated together at the single operating point $(x_0,u_0)$ where
$f(x_0,u_0)=0$. Every entry in the ``s-domain relations'' table you will build
up over the unit is, underneath, one of these two partial derivatives.
\end{keyidea}

% ====================================================================
\section*{Question 4 --- Gearing with load friction retained \athome}

Redo the gear-reflection idea of Question~2, but for a single gear stage where
the load's own friction is \emph{not} negligible. A motor with
$J_m=\SI{0.05}{\kilogram\metre\squared}$ and
$b=\SI{0.06}{\newton\metre\second\per\radian}$ (the same motor as
Question~1) drives, through a single gear pair of ratio
$n=\theta_L/\theta_m=1/5$, a load of inertia
$J_L=\SI{2.5}{\kilogram\metre\squared}$ and \emph{load-side} viscous friction
$b_L=\SI{0.4}{\newton\metre\second\per\radian}$.

\begin{enumerate}[label=(\alph*)]
\item Find the equivalent inertia $J_{\text{eq}}=J_m+n^2J_L$ and the
      equivalent damping $b_{\text{eq}}=b+n^2b_L$ at the motor shaft.
\item Using the purely mechanical (no electrical circuit) transfer function
      \[
        \frac{\theta_m(s)}{T_m(s)}=\frac{1}{s(J_{\text{eq}}s+b_{\text{eq}})},
      \]
      find the time constant $J_{\text{eq}}/b_{\text{eq}}$.
\item By what percentage does ignoring $b_L$ entirely (as Question~2 did)
      change this time constant? Is it a bigger or smaller effect than
      ignoring $L_a$ was in Question~1?
\end{enumerate}

\begin{solution}
\textbf{(a)} With $n^2=1/25=0.04$:
\[
  J_{\text{eq}}=0.05+0.04(2.5)=0.05+0.1=\SI{0.15}{\kilogram\metre\squared},
\]
\[
  b_{\text{eq}}=0.06+0.04(0.4)=0.06+0.016=\SI{0.076}{\newton\metre\second\per\radian}.
\]

\medskip
\textbf{(b)}
\[
  \tau=\frac{J_{\text{eq}}}{b_{\text{eq}}}=\frac{0.15}{0.076}=\SI{1.974}{\second}.
\]

\medskip
\textbf{(c)} Ignoring $b_L$ means using $b_{\text{eq}}\to b=0.06$ instead of
$0.076$, giving $\tau\to0.15/0.06=\SI{2.5}{\second}$ --- a change of
\[
  \frac{2.5-1.974}{1.974}\times100\%\approx27\%,
\]
substantially bigger than the $4.4\%$ effect of ignoring $L_a$ in
Question~1. Load friction reflected through even a modest reduction gear is
not automatically negligible, and should be checked rather than assumed away
by habit.
\end{solution}

% ====================================================================
\section*{Question 5 --- Linearizing quadratic (turbulent) drag \athome}

A rotational load of inertia $J=\SI{0.8}{\kilogram\metre\squared}$ is driven by
torque $T(t)$ against turbulent aerodynamic drag, which for this load is
proportional to the \emph{square} of the speed rather than to the speed
itself:
\[
  J\dot\omega(t) = T(t) - c\,\omega(t)\,|\omega(t)|,
  \qquad \omega=\dot\theta,\quad c=\SI{0.02}{\newton\metre\second\squared\per\radian\squared}.
\]
(The $|\omega|$ keeps the drag opposing the motion whichever way the load
turns; for $\omega>0$ it is simply $c\omega^2$.)

\begin{enumerate}[label=(\alph*)]
\item Suppose the load is to be held at a constant speed $\omega_0=\SI{10}{\radian\per\second}$
      ($\omega_0>0$). Find the torque $T_0$ needed at this operating point.
\item Linearize the drag torque about $\omega_0$ to find an \emph{effective
      linear} damping coefficient $b_{\text{eff}}$, defined by
      $\Delta T_{\text{drag}}\approx b_{\text{eff}}\Delta\omega$.
\item Hence write the linearized equation for $\Delta\omega$ in terms of
      $\Delta T$, and find its time constant.
\item A common mistake is to linearize by simply evaluating the \emph{drag
      itself} at $\omega_0$, i.e.\ to treat $b_{\text{eff}}=c\omega_0$ rather
      than using the derivative. Compute that (wrong) value and compare it
      with your answer to (b). Which is bigger, and by what factor?
\end{enumerate}

\begin{solution}
\textbf{(a)} At equilibrium $\dot\omega=0$, so
$T_0=c\omega_0^2=(0.02)(10)^2=\SI{2}{\newton\metre}$.

\medskip
\textbf{(b)} For $\omega>0$ the drag is $c\omega^2$, so
\[
  b_{\text{eff}}=\left.\frac{d}{d\omega}\big(c\omega^2\big)\right|_{\omega_0}
  =2c\omega_0=2(0.02)(10)=\SI{0.4}{\newton\metre\second\per\radian}.
\]

\medskip
\textbf{(c)} $J\Delta\dot\omega=\Delta T-b_{\text{eff}}\Delta\omega$, i.e.
\[
  \boxed{\;0.8\,\Delta\dot\omega+0.4\,\Delta\omega=\Delta T\;},
  \qquad \tau=\frac{J}{b_{\text{eff}}}=\frac{0.8}{0.4}=\SI{2}{\second}.
\]

\medskip
\textbf{(d)} The naive value is $c\omega_0=(0.02)(10)=\SI{0.2}{\newton\metre\second\per\radian}$
--- exactly \emph{half} of the correct $b_{\text{eff}}=0.4$ found by
differentiating. This is not a coincidence: for any power-law drag
$c\omega^{p}$, the correct linearized coefficient is $pc\omega_0^{p-1}$, which
for $p=2$ is $2c\omega_0$, twice the naive $c\omega_0$.
\end{solution}

\begin{pitfall}
Linearization is differentiation, not substitution. Evaluating the nonlinear
function \emph{at} the operating point gives you the value $T_0$ there (useful
for part (a)); it does not give you the \emph{slope}, which is what governs
the response to a small deviation and is what you actually need for the linear
model. This is the single most common error in linearization problems,
and Question~5(d) is designed to make it show up numerically rather than
staying a rule you forget under pressure.
\end{pitfall}

% ====================================================================
\vfill
\begin{center}
\rule{0.5\textwidth}{0.4pt}\\[4pt]
{\small\itshape
\ifsolutions
  End of solutions. Week~3: transfer functions and the $s$-plane.
\else
  Solutions to the homework questions are released with the Week~3 sheet.\\
  Next week: transfer functions and the $s$-plane.
\fi}
\end{center}

\end{document}
